% =====================================================================
%  Chapter: Reduced-order hydrodynamic model and gait optimisation (MuJoCo)
%  Source code: swimopt/ ; data behind every figure: swimopt/results_thesis/
% =====================================================================
\providecommand{\chCFD}{the CFD chapter}
% Screenshot helper: shows the image if figures/<file> exists, otherwise a framed placeholder.
\providecommand{\shot}[3]{\IfFileExists{figures/#1}{\includegraphics[width=#2]{#1}}%
  {\fbox{\parbox{\dimexpr#2-2\fboxsep-2\fboxrule\relax}{\centering\small\vspace*{9mm}\textbf{MuJoCo screenshot \texttt{\detokenize{#1}}}\\[3pt]#3\par\vspace*{9mm}}}}}

\chapter{Reduced-Order Hydrodynamic Model and Trajectory Optimisation}
\label{ch:mujoco}

The CFD studies in \chCFD{} resolve the flow around a single leg or the bare hull
with high fidelity, but one unsteady case takes hours to days. Finding a good
paddling trajectory needs the opposite trade-off: thousands of whole-robot
evaluations, each of which only has to rank candidate gaits correctly. This chapter
develops a reduced-order model for that purpose. A quasi-steady, per-link
hydrodynamic force model is attached to a rigid-body simulation in MuJoCo
\cite{Todorov2012}, the joint trajectories are written as a low-dimensional
parameter vector, and an evolution strategy (CMA-ES) searches that vector. The
same pipeline is applied to two robots: an open toy quadruped with drag-type
flippers, and the BODY2 robot with passively pitching lift-type flukes.

The chapter has four aims:
\begin{enumerate}
  \item define the hydrodynamic and gait models precisely enough to reproduce them (\cref{sec:mj-hydro,sec:mj-gait});
  \item verify the implementation and the evaluation protocol, including three
        errors that were found and corrected during this work (\cref{sec:mj-verify});
  \item show what the optimiser finds for a drag-type robot, and why its answer
        led to the lift-type fluke design (\cref{sec:mj-toy});
  \item quantify the fluke design space of BODY2 with the uncalibrated model and
        state what must be calibrated before these numbers can be trusted
        (\cref{sec:mj-fluke,sec:mj-discussion}).
\end{enumerate}
All speeds in this chapter are model predictions with coefficients taken from the
literature or from the hull CFD. None of them is a measured swimming speed.

% ---------------------------------------------------------------------
\section{Simulation platform and pipeline}
\label{sec:mj-pipeline}

\paragraph{Platform choice.}
The first implementation used Webots. BODY2's legs are closed parallelogram and
four-bar linkages. URDF cannot describe closed chains, so the loops had to be
closed by hand after import. In Webots this produced two problems that could not be
removed by parameter tuning. First, a closed four-bar has two assembly modes, and
with all joints initialised at zero the solver placed different legs on different
branches (``one leg always behaves differently''). Second, several robot instances
in one world influenced each other's assembly, so parallel comparison of gaits was
unreliable. MuJoCo \cite{Todorov2012} was adopted instead. It simulates in
generalised coordinates, closes loops with explicit equality constraints, and resets
a state deterministically by overwriting \texttt{qpos}/\texttt{qvel}; identical
parameters give bit-identical results (\cref{sec:mj-verify}). MuJoCo has no free
surface and no built-in fluid model suitable for partially immersed bodies. Both
were therefore implemented as an external force module
(\cref{sec:mj-hydro}), which also keeps every modelling assumption explicit.

\paragraph{Pipeline.}
The software (\texttt{swimopt}) is split into four layers so that a new robot
changes only the first (\cref{fig:mj-pipeline}). The optimiser sees only a function
$J(\bm x)$ of a normalised parameter vector $\bm x\in[0,1]^n$.

\begin{figure}[t]
\centering
\begin{tikzpicture}[
  font=\footnotesize, node distance=4mm and 5mm,
  box/.style={draw, rounded corners=2pt, align=center, minimum height=9mm, text width=27mm, fill=#1!10},
  box/.default=gray, arr/.style={-{Latex[length=2mm]}, thick}]
  \node[box=blue] (model) {\textbf{\textcircled{1} Model}\\ \texttt{robots/*.xml}\\ MJCF, from URDF};
  \node[box=teal, right=of model] (hydro) {\textbf{\textcircled{2} Physics}\\ Morison elements\\ \texttt{hydro.py}};
  \node[box=teal, below=of hydro] (foil) {\textbf{\textcircled{2}b Physics}\\ quasi-steady foil\\ \texttt{hydro\_foil.py}};
  \node[box=orange, right=of hydro] (gait) {\textbf{\textcircled{3} Control}\\ Fourier gait + duty\\ \texttt{gait.py}};
  \node[box=gray, below=of gait] (sim) {\textbf{rollout}\\ reset, settle, run\\ \texttt{simulate.py}};
  \node[box=red, right=of gait] (opt) {\textbf{\textcircled{4} Optimiser}\\ CMA-ES\\ \texttt{optimize.py}};
  \node[box=gray, below=of opt] (fit) {\textbf{score} $J(\bm x)$\\ speed $-$ penalties\\ feasibility checks};
  \node[box=blue, below=of model] (imp) {\texttt{import\_model.py}\\ meshes, actuators};
  \node[box=gray, below=of fit, yshift=-1mm] (ik) {\texttt{leg\_v3.py}: four-bar IK\\ \texttt{export\_fluke\_gait.py}};
  \draw[arr] (imp) -- (model);
  \draw[arr] (model) -- (hydro);
  \draw[arr] (hydro) -- (gait);
  \draw[arr] (foil.east) -- (sim.west);
  \draw[arr] (hydro) -- (foil);
  \draw[arr] (gait) -- (sim);
  \draw[arr] (sim) -- (fit);
  \draw[arr] (fit) -- (opt);
  \draw[arr] (opt.north) -- ++(0,5mm) -| node[pos=0.25, above]{candidate $\bm x$} (gait.north);
  \draw[arr, dashed] (ik) -- (fit);
\end{tikzpicture}
\caption{Structure of the \texttt{swimopt} pipeline. Only layer \textcircled{1} depends on
the robot; hydrodynamic coefficients are assigned to links by name rules in a
configuration file. The four-bar inverse kinematics of BODY2 enter only as
feasibility checks and in the firmware export.}
\label{fig:mj-pipeline}
\end{figure}

% ---------------------------------------------------------------------
\section{Hydrodynamic model}
\label{sec:mj-hydro}

Two element types are used. Hull, links and drag-type paddles are \emph{Morison
elements}: forces depend on the local velocity of each collision geometry. The BODY2
flukes are \emph{foil elements}: forces depend on angle of attack and follow a
quasi-steady wing model. Both are applied as external wrenches
(\texttt{xfrc\_applied}) before every integration step.

\subsection{Morison elements}
\label{sec:mj-morison}

For every immersed geometry $i$ with volume $V_i$, orientation $\bm R_i$ and
frontal areas $A_{i,k}$ normal to its local axes $k\in\{x,y,z\}$, the applied force is
\begin{equation}
  \bm F_i \;=\; \underbrace{\rho g V_i f_i\,\bm e_z}_{\text{buoyancy}}
  \;-\;\underbrace{\tfrac12\rho\,\bm R_i\big(\bm C_{d,i}\odot\bm A_i\odot\bm v^{\ell}_i\odot|\bm v^{\ell}_i|\big)f_i}_{\text{anisotropic quadratic drag}}
  \;-\;\underbrace{c_{v,i}\,\bm v_i f_i}_{\text{linear damping}}
  \;+\;\underbrace{m_{a,i}\,g\,\bm e_z}_{\text{added-mass compensation}},
  \label{eq:mj-morison}
\end{equation}
where $\bm v_i$ is the world velocity of the geometry centre, $\bm v^{\ell}_i=\bm R_i^{\top}\bm v_i$
its components in the body-fixed frame, $\odot$ the element-wise product, and
$f_i\in[0,1]$ the immersed fraction. The force acts at the geometry centre, so the
wrench on the parent body also contains the moment
$(\bm p_i-\bm c_b)\times\bm F_i$ about the body's centre of mass $\bm c_b$. This is the
quasi-steady drag and inertia decomposition of Morison et al.\ \cite{Morison1950},
applied per link as is common for underwater vehicles \cite{Fossen1994} and used for
paddling robots by Chen et al.\ \cite{Chen2021,Chen2023}.

\paragraph{Anisotropy.}
Each element carries three drag coefficients. For a thin paddle the coefficient
normal to its face is large and the two edge-wise coefficients are small
(\cref{tab:mj-params}). This asymmetry is what allows a reciprocating paddle to produce
net thrust at all; with an isotropic coefficient the force would always be exactly
opposite to the paddle velocity.

\paragraph{Partial immersion.}
The immersed fraction uses the true vertical half-extent $h_i$ of the geometry at
its current orientation, for example
$h_i=\sum_k|R_{i,zk}|\,s_{i,k}$ for a box with half-sizes $s_{i,k}$ and
$h_i=|R_{i,zz}|\,\ell_i+r_i$ for a capsule. With the water surface at $z_w$,
\begin{equation}
  f_i=\operatorname{clip}\!\left(\frac{z_w-(p_{i,z}-h_i)}{2h_i},\,0,\,1\right).
  \label{eq:mj-immersion}
\end{equation}
The linear profile is exact for the volume of an upright box and an approximation
otherwise. Mesh geometries are replaced by their vertex bounding box scaled by a fill
factor (default $0.55$).

\paragraph{Added mass.}
An explicit term $-m_a\dot{\bm v}$ uses the acceleration of the current step and
is numerically stiff for light links. Instead, the added mass
$m_{a,i}=C_{a,i}\rho V_i$ is added to the inertial mass of the parent body once at
initialisation, with the rotational inertia scaled by the same factor, and its weight
is cancelled by the constant upward force in \cref{eq:mj-morison}. The body then has
the correct weight and an increased inertia. This is unconditionally stable, but it
makes the added mass isotropic and independent of immersion, whereas the physical
added mass of a plate is large only normal to its face \cite{Newman1977}. The
consequence is discussed in \cref{sec:mj-discussion}.

\subsection{Quasi-steady foil elements for the passive fluke}
\label{sec:mj-foil}

The BODY2 fluke is a thin plate on a pin at the end of the shank. It rotates freely
about the pin against a leaf spring. Its thrust comes from the force normal to the
relative flow, so a model based on angle of attack is needed. Each foil is a box
with chord $c$ along its local $x$ axis, span $b$ along $y$ and normal $\hat{\bm z}$;
$S=cb$ and $\mathit{AR}=b/c$ (\cref{fig:mj-foil-sketch}).

\begin{figure}[t]
\centering
\begin{tikzpicture}[scale=1.25, font=\small, >={Latex[length=2mm]}]
  % shank
  \draw[line width=1.6pt, gray!70] (-2.6,1.0) -- (0,0) node[pos=0.1, above, black]{shank};
  \fill (0,0) circle (1.6pt) node[above left]{pin};
  % spring symbol
  \draw[decorate, decoration={coil, aspect=0.6, segment length=2.2mm, amplitude=1.4mm}] (0,0) ++(160:0.45) arc (160:215:0.45);
  \node at (-0.72,-0.35) {$k$};
  % foil plate rotated by theta
  \begin{scope}[rotate=-18]
    \draw[fill=black!75] (0.08,-0.05) rectangle (2.4,0.05);
    \draw[dashed, gray] (0,0) -- (2.9,0) node[right, black]{$\hat{\bm x}$ chord};
    \fill[blue] (0.08+0.25*2.32,0) circle (1.4pt) node[above=2pt, blue]{$\tfrac14 c$};
    \fill[red] (0.08+0.75*2.32,0) circle (1.4pt) coordinate (p34) node[above=2pt, red]{$\tfrac34 c$};
    \draw[->, thick] (0.08+0.25*2.32,0) -- ++(0.35,1.15) node[above]{$\bm L$};
    \draw[->, thick] (0.08+0.25*2.32,0) -- ++(0.55,-0.17) node[right]{$\bm D$};
  \end{scope}
  % relative flow at 3/4 chord
  \draw[->, thick, teal] ($(p34)+(-1.0,-0.62)$) -- (p34) node[pos=0.1, below, teal]{relative flow $\bm w=-\bm v_{3/4}$};
  \draw[gray] (p34) ++(198:0.55) arc (198:212:0.55);
  \node[font=\footnotesize] at ($(p34)+(-0.85,-0.14)$) {$\alpha$};
  \node[align=left, font=\footnotesize, anchor=west] at (3.3,0.9) {$\theta$: passive pitch about the pin\\ stop at $\pm40^\circ$};
\end{tikzpicture}
\caption{Foil element of the BODY2 fluke. The relative flow is taken at the
three-quarter-chord point, the aerodynamic force acts at the quarter-chord point, and
the pitching moment about the pin that drives the passive rotation $\theta$ follows
from the lever arm between the two.}
\label{fig:mj-foil-sketch}
\end{figure}

Following the quasi-steady reduction of Theodorsen's theory \cite{Theodorsen1935},
the relative flow is evaluated at the three-quarter-chord point, where the
translational and rotational velocity contributions combine:
$\bm w=-\bm v_{3/4}$, with $w_x=\bm w\cdot\hat{\bm x}$ and $w_z=\bm w\cdot\hat{\bm z}$.
The angle of attack and dynamic pressure are
\begin{equation}
  \alpha=\operatorname{atan2}(w_z,w_x),\qquad q=\tfrac12\rho\,(w_x^2+w_z^2)\,S\,f .
\end{equation}
Lift acts normal to the in-plane flow and drag along it,
$\bm F=q\,(C_L\,\hat{\bm l}+C_D\,\hat{\bm w})$, applied at the quarter-chord point. The
coefficients blend attached-flow thin-wing values with a post-stall flat-plate model:
\begin{align}
  C_L(\alpha)&=\big[\sigma\,C_{L\alpha}+(1-\sigma)\,C_N\big]\sin\alpha\cos\alpha, \\
  C_D(\alpha)&=\sigma\Big(C_{D0}+\frac{(C_{L\alpha}\sin\alpha\cos\alpha)^2}{\pi e\,\mathit{AR}}\Big)+(1-\sigma)\,C_N\sin^2\alpha,\\
  \sigma(\alpha)&=\operatorname{clip}\!\left(\frac{1.6\,\alpha_s-|\alpha|}{0.6\,\alpha_s},0,1\right),
  \qquad C_{L\alpha}=\frac{2\pi\,\mathit{AR}}{\mathit{AR}+2}.
\end{align}
The lift slope is Helmbold's low-aspect-ratio correction \cite{Helmbold1942}. The
weight $\sigma$ is one for $|\alpha|\le\alpha_s$ and zero beyond $1.6\,\alpha_s$
(\cref{fig:mj-foil-coeffs}). Buoyancy and added mass
($m_a=C_a\rho\pi(c/2)^2b$) are handled as for Morison elements. The passive hinge is a
MuJoCo joint with the leaf-spring stiffness $k$, a small damping, and a
$\pm\SI{40}{\degree}$ range limit.

\begin{figure}[t]
  \centering
  \includegraphics[width=\linewidth]{fig_foil_coeffs}
  \caption{Foil coefficients for the BODY2 fluke ($\mathit{AR}=3.13$,
  $C_{L\alpha}=\SI{3.83}{\per\radian}$, $\alpha_s=\SI{15}{\degree}$, $C_N=1.8$). Left:
  lift and drag coefficients; the grey bands are the blending ranges. Right:
  attached-flow weight $\sigma$.}
  \label{fig:mj-foil-coeffs}
\end{figure}

\begin{table}[t]
\centering
\caption{Hydrodynamic parameters used in this chapter. None has been calibrated
against the physical robot yet, except the hull's surge coefficients, which were
fitted to the bare-hull CFD at $0.1$ and $\SI{0.2}{\metre\per\second}$ (\chCFD).}
\label{tab:mj-params}
\footnotesize
\setlength{\tabcolsep}{4pt}
\begin{tabular}{@{}llccc@{}}
\toprule
robot & element (name rule) & $C_d$ $(x,y,z)$ & $C_a$ & $c_v$ [\si{\newton\second\per\metre}] \\
\midrule
toy     & flipper (\texttt{flip})   & $(2.2,\;0.10,\;0.10)$ & 1.0  & 0.05 \\
        & leg link (\texttt{link})  & $(0.5,\;0.50,\;0.25)$ & 0.3  & 0.02 \\
        & trunk (\texttt{trunk})    & $(0.9,\;1.60,\;1.80)$ & 0.2  & 0.10 \\
\midrule
BODY2   & thigh, shank (\texttt{link}) & $(0.6,\;0.6,\;0.25)$ & 0.3 & 0.02 \\
fluke   & hull (\texttt{trunk})     & $(0.093,\;1.2,\;1.5)$ & 0.15 & 0.0154 \\
\cmidrule(l){2-5}
        & fluke foil (\texttt{foil}) & \multicolumn{3}{l}{$c=\SI{19.2}{mm}$, $b=\SI{60}{mm}$, $\alpha_s=\SI{15}{\degree}$, $C_{D0}=0.02$,} \\
        & & \multicolumn{3}{l}{$e=0.9$, $C_N=1.8$, $C_a=1.0$} \\
        & passive hinge & \multicolumn{3}{l}{$k=\SI{2.5}{mN.m/rad}$, range $\pm\SI{40}{\degree}$} \\
\bottomrule
\end{tabular}
\end{table}

% ---------------------------------------------------------------------
\section{Gait parameterisation and optimisation}
\label{sec:mj-gait}

\subsection{Truncated Fourier series with duty-cycle warp}
Every position actuator $i$ follows
\begin{equation}
  q_i(t)=o_i+r(t)\sum_{h=1}^{H}A_{i,h}\sin\!\big(h\,\vartheta(t)+\varphi_{i,h}\big),
  \label{eq:mj-gait}
\end{equation}
with offset $o_i$, amplitudes $A_{i,h}$, phases $\varphi_{i,h}$ and a start-up ramp
$r(t)=\min(1,t/\SI{1.5}{s})$. The phase variable $\vartheta$ is warped so that a
fraction $d$ of the period (the power-stroke duty cycle) covers the first half-wave:
with $\tau=(ft)\bmod 1$,
\begin{equation}
  \vartheta=\begin{cases}\pi\,\tau/d, & \tau<d,\\[2pt] \pi+\pi\,(\tau-d)/(1-d), & \tau\ge d.\end{cases}
\end{equation}
For $d=0.5$ this is the ordinary sinusoid. Actuators can be grouped so that they share
amplitude and offset (for example all hip joints), while phases stay individual. With
$n_u$ actuators, $n_A$ amplitude groups and $n_o$ offset groups, the parameter
vector has
\begin{equation}
  n = 2 + H\,(n_A+n_u) + n_o
\end{equation}
entries (frequency, duty, amplitudes, phases, offsets), each mapped linearly from $[0,1]$
to a physical range. For the toy robot ($n_u=12$, three shared groups, $H=2$) this
gives $n=35$; for BODY2 ($n_u=4$, one group, $H=1$, duty fixed) it gives $n=8$, of
which seven are optimised.

\subsection{Why the flipper needs a second harmonic}
\label{sec:mj-harmonic}
For $H=1$ every joint satisfies $q_i(t+T/2)-o_i=-(q_i(t)-o_i)$: the second half
of each cycle is the first half reflected about the offset posture. For the toy robot,
whose legs hang vertically and are fore--aft symmetric about their hips at zero
offset, this reflects each leg's stroke about the vertical, so quasi-steady forces from
the two half-cycles cancel. A real paddle breaks the symmetry by feathering: it
opens on the power stroke and closes on the recovery stroke, which is an effect at
twice the stroke frequency. The term $A_{i,2}\sin(2\vartheta+\varphi_{i,2})$ is
unchanged by the half-period shift and represents exactly this.

\Cref{fig:mj-harmonic} tests the argument on the toy robot with a fixed diagonal
gait (hip, knee and flipper first harmonics fixed, zero offsets), sweeping only the
flipper's second-harmonic amplitude $A_2$ and phase $\psi$. Without the second
harmonic the speed is below \SI{0.2}{mm/s} (numerical noise of the settling
transient). With it, the speed follows a sinusoid in $\psi$ whose amplitude grows
with $A_2$ up to \SI{9.1}{cm/s}. Shifting $\psi$ by \SI{180}{\degree} reverses the
speed exactly (to $10^{-13}$\,\si{mm/s}): the phase of the feathering relative to the
stroke decides the swimming direction. A first-harmonic-only parameterisation would
therefore search a space that contains no thrust at all for this robot.

\begin{figure}[t]
  \centering
  \includegraphics[width=\linewidth]{fig_harmonic_map}
  \caption{Toy robot, diagonal gait with all first harmonics fixed: forward speed as a
  function of the flipper's second-harmonic amplitude $A_2$ and phase $\psi$. Left: map;
  right: sections. The row $A_2=0$ is zero everywhere, and $v(\psi+\SI{180}{\degree})=-v(\psi)$.}
  \label{fig:mj-harmonic}
\end{figure}

\paragraph{Structured reduction.}
In the 35-dimensional toy space most phase combinations are not recognisable gaits. A
reduced space fixes the inter-leg first-harmonic phases to a named pattern (diagonal,
front--back, left--right, wave, in-phase) and couples the second-harmonic phases by
$\varphi_{i,2}=\psi+2\varphi_{i,1}$. The coupling gives every leg the same feathering
timing relative to its own stroke. The remaining search variables are frequency, duty,
three first- and three second-harmonic group amplitudes, $\psi$ and three offsets:
12 dimensions.

\subsection{Objective and feasibility}
One evaluation (``rollout'') resets the robot, lets it settle in its initial posture for
$t_s$ (\SIrange{1.0}{1.5}{s}), and then runs the gait for $T=\SI{8}{s}$ with
$\Delta t=\SI{2}{ms}$. The score is
\begin{equation}
  J(\bm x)=\bar u-w_\psi\frac{|\Delta\psi|}{T}-w_E\bar P-w_\theta\,\mathrm{std}(\theta_{\text{pitch}})
  -w_z\,\mathrm{std}(z)-w_{\dot q}\,\Big(\frac{\max|\dot q_{\text{cmd}}|}{\dot q_{\max}}-1\Big)^{+},
  \label{eq:mj-objective}
\end{equation}
where $\bar u$ is the displacement along the initial heading divided by $T$,
$\Delta\psi$ the net yaw change, $\bar P$ the mean absolute actuator power, and the
standard deviations of trunk pitch and heave are taken over the second half of the run.
The fitness is a ruler, not a result: its value depends on the weights and is not
compared across settings. Only physical quantities are reported.

For BODY2 the objective includes $w_\theta=0.5$, $w_z=2.0$, $w_\psi=0.3$ and a servo
speed limit $\dot q_{\max}=\SI{14}{rad/s}$. Before a rollout, both extreme shank angles
$q_r=o\pm A$ are checked against the four-bar inverse kinematics of the real leg
(transmission angle $\ge\SI{10}{\degree}$, knee angle $\ge\SI{45}{\degree}$) and against
the fluke--hull clearance limit $q_r\ge\SI{15}{\degree}$. Infeasible candidates score
$-1$ without being simulated. Rollouts whose trunk speed exceeds \SI{5}{m/s} or whose
joint speeds exceed \SI{100}{rad/s} are aborted and score $-1000$.

\subsection{CMA-ES}
The covariance matrix adaptation evolution strategy \cite{Hansen2001,Hansen2016}
samples $\lambda$ candidates $\bm x_k\sim\bm m+\sigma\,\mathcal N(\bm 0,\bm C)$, ranks them
by $J$, moves the mean $\bm m$ to a weighted average of the best $\mu$, and adapts the
step size $\sigma$ (cumulative step-size adaptation) and the covariance $\bm C$ (rank-one
and rank-$\mu$ updates). Only ranks are used, so the scale of $J$ does not matter, and
no gradient is needed. The latter is essential here because contact with joint limits,
feasibility checks and the divergence guard make $J$ discontinuous. Box bounds
$[0,1]^n$ are handled by the \texttt{cma} package. Settings: $\lambda=12$ for the full
toy space and BODY2, $\lambda=10$ for the 12-D reduced space, $\sigma_0=0.25$ (toy) and
$0.3$ (BODY2). A comparable CMA-ES workflow for gait selection of a fish-like robot
is reported in \cite{FishCMAES2021}. The overall workflow (parameterise, simulate,
optimise) follows Lee et al.\ \cite{Lee2025}. They used a differentiable simulator and a
gradient method; MuJoCo with the forces above is not differentiable, so a
derivative-free optimiser is used.

% ---------------------------------------------------------------------
\section{Implementation and verification}
\label{sec:mj-verify}

\paragraph{Numerical checks.}
\Cref{tab:mj-verify} summarises the checks. The vectorised force routine agrees with a
per-geometry reference loop to \SI{1.8e-15}{N} over 1500 steps of a moving gait and is
4.8 times faster (\SI{0.15}{ms} vs.\ \SI{0.73}{ms} per step for 13 elements). This
comparison exposed an error in the reference loop itself: a NumPy integer compared
with a MuJoCo enum inside a Python \texttt{in} test evaluates false, so capsules were
treated as spheres. The production path was not affected; the reference was corrected
before the comparison. A second vectorised implementation of the foil routine, used for
the parameter sweeps, agrees with the original to $10^{-14}$ in all reported metrics.

\paragraph{Flotation.}
With the legs held still, the time-averaged buoyancy after settling equals the weight
to within \SI{0.06}{\percent} (toy) and \SI{0.03}{\percent} (BODY2)
(\cref{fig:mj-checks}, left). The heave oscillation is lightly damped and still has a
peak-to-peak amplitude of about \SI{3}{mm} after \SI{8}{s}. BODY2 settles with its
reference point O2 at \SI{-6.4}{mm} and a bow-up trim of \SI{1.3}{\degree}. The target
waterline from the hull design is \SI{-11}{mm}, so the provisional ballast
(\SI{0.85}{kg}) is too light. This is one of the calibration items in
\cref{sec:mj-discussion}.

\paragraph{Virtual coast-down.}
The robot is launched at \SI{0.25}{m/s} with the legs held at \SI{40}{\degree}, and
$m_{\text{eff}}\dot u=-a u^2-b u$ is fitted to the simulated deceleration
(\cref{fig:mj-checks}, right). The fit recovers $a=\SI{1.93}{N.s^2/m^2}$ and
$b=\SI{0.144}{N.s/m}$ with $m_{\text{eff}}=\SI{1.65}{kg}$ (including added mass). The
hull's own coefficients, weighted by its immersed fraction ($0.91$), are $a=0.24$ and
$b=0.014$. In the model the appendages, above all the flukes held at a large angle,
therefore account for about \SI{87}{\percent} of the quadratic and \SI{90}{\percent} of the
linear surge resistance at rest posture. Hull CFD alone would
underestimate the robot's drag, and the physical coast-down test should be done with
the legs attached.

\paragraph{Evaluation protocol.}
Three errors in the evaluation protocol were found during this work. All results in
this chapter were computed after correcting them.
\begin{enumerate}
  \item \emph{Start-up step.} The settling phase commanded zero joint angles, and the
    gait started from its offset angles. The stiff toy servos then moved the light
    joints at more than \SI{100}{rad/s} in the first step, which triggered the
    divergence guard. Of 24 random candidates, 16 were rejected for this reason. After
    commanding the offset posture during settling, none were.
  \item \emph{Settling stroke.} Even with a slow servo, moving the legs from zero to
    the offset posture during settling is itself a stroke. On BODY2 the resulting glide
    persisted into the scored interval. With zero stroke amplitude the robot still
    covered \SI{1.6}{cm/s}. The joints are now placed directly in the start posture at
    reset. The same zero-amplitude \emph{null test} then gives \SI{-0.12}{mm/s}
    (BODY2) and \SI{-3.5}{mm/s} (toy).
  \item \emph{Pitch metric.} The pitch penalty originally used the RMS of the pitch
    angle, which is dominated by the static trim (\SI{1.3}{\degree}) rather than the
    nodding it was meant to measure (\SIrange{0.1}{0.4}{\degree}). It now uses the
    standard deviation, like the heave term.
\end{enumerate}
Errors 1 and 2 had affected all earlier results. The earlier BODY2 optimum of
\SI{3.2}{cm/s} contained about \SI{1.6}{cm/s} of settling glide. The earlier toy result
had in addition been obtained with joint ranges interpreted in degrees (MJCF default)
instead of radians.

\begin{table}[t]
\centering
\caption{Verification checks of the implementation.}
\label{tab:mj-verify}
\small
\begin{tabularx}{\linewidth}{@{}lXl@{}}
\toprule
check & criterion & result \\
\midrule
vectorised vs.\ reference forces & max.\ difference over 1500 steps & \SI{1.8e-15}{N} \\
determinism & two rollouts, same $\bm x$ & bit-identical \\
buoyancy balance & mean buoyancy / weight, legs still & $1.0006$ (toy), $1.0003$ (BODY2) \\
null test & speed at zero stroke amplitude & \SI{-3.5}{mm/s} (toy), \SI{-0.12}{mm/s} (BODY2) \\
second-harmonic symmetry & $v(\psi)+v(\psi+\SI{180}{\degree})$ & $<10^{-13}$\,\si{mm/s} \\
cost per rollout & \SI{9}{s} simulated, one CPU core & $\approx\SI{1.2}{s}$ (toy), $\approx\SI{2.2}{s}$ (BODY2) \\
\bottomrule
\end{tabularx}
\end{table}

\begin{figure}[t]
  \centering
  \includegraphics[width=\linewidth]{fig_model_checks}
  \caption{Left: settling from the design height with the legs held still. Right:
  simulated coast-down of BODY2 from \SI{0.25}{m/s}, the fitted quadratic-plus-linear
  resistance, and the deceleration that the hull coefficients alone would give.}
  \label{fig:mj-checks}
\end{figure}

% ---------------------------------------------------------------------
\section{Case A: toy quadruped with drag-type flippers}
\label{sec:mj-toy}

\begin{figure}[t]
  \centering
  \shot{shot_toy_quad.png}{0.48\linewidth}{\texttt{python view.py config.json --demo}\\ with \texttt{"model": "robots/toy\_quad.xml"};\\ oblique side view, water plane visible}\hfill
  \shot{shot_body2_fluke.png}{0.48\linewidth}{\texttt{python view.py config\_fluke.json}\\ close-up of one leg: fixed thigh,\\ shank, pin and fluke}
  \caption{The two models. Left: toy quadruped, three pitch joints per leg (hip, knee,
  flipper wrist), \SI{0.78}{kg}. Right: BODY2 with passive-pitch flukes, one actuated
  shank per leg, \SI{1.35}{kg} including ballast.}
  \label{fig:mj-models}
\end{figure}

The toy robot (\cref{fig:mj-models}, left) has a $200\times100\times\SI{50}{mm}$ trunk
and four legs with hip, knee and wrist joints driven by position servos
($k_p=\SI{60}{N.m/rad}$, $|\tau|\le\SI{12}{N.m}$), and a
$52\times\SI{36}{mm}$ flipper on each wrist. These actuators are far stronger than any
servo of a robot this size. The toy is therefore used for methodological questions
only: which parameterisation, which search space, and which propulsion mode the optimiser
chooses. Its speeds are not a prediction for BODY2. All toy runs use $w_\psi=1.0$ and no
energy, pitch or heave penalty.

\subsection{Search space: full versus structured}
\Cref{fig:mj-convergence} compares CMA-ES on the full 35-dimensional space (two seeds,
480 rollouts each) with the structured 12-dimensional space (two seeds per pattern, 150
rollouts each). In the full space the best speeds were \SI{12.5}{cm/s} and
\SI{4.9}{cm/s} for the two seeds. Both incumbents were found within 72 rollouts and never
improved afterwards, and the step size did not contract ($\sigma\approx0.18$--$0.25$ at the
end). The search had not converged, and the two seeds disagree by a factor of 2.5.
The structured space reached \SIrange{20.2}{22.2}{cm/s} with the diagonal pattern and
\SIrange{14.2}{26.8}{cm/s} with the front--back pattern within 150 rollouts. Fixing the inter-leg phases and tying the feathering phase
to the stroke removes directions that only rotate the robot or cancel thrust between legs,
and this helps the optimiser more than the extra freedom of the full space.

\begin{figure}[t]
  \centering
  \includegraphics[width=\linewidth]{fig_convergence}
  \caption{Toy robot. Left: speed of the best candidate found so far. Right: CMA-ES step
  size. Grey: full 35-D space; blue, red: structured 12-D space with diagonal and
  front--back phase patterns. Solid and dashed lines are two random seeds.}
  \label{fig:mj-convergence}
\end{figure}

\subsection{Gait patterns}
\Cref{fig:mj-gait-patterns} lists the best result per pattern. Diagonal
(\SIrange{20.2}{22.2}{cm/s}) and front--back (\SIrange{14.2}{26.8}{cm/s}) are the
fastest, followed by in-phase (\SIrange{15.0}{16.4}{cm/s}), wave
(\SIrange{5.8}{9.8}{cm/s}) and left--right (\SIrange{2.0}{8.6}{cm/s}). The
seed-to-seed spread is as large as some differences between patterns, so with two seeds
only the grouping ``diagonal and front--back fast, wave and left--right slow'' is
supported. All of these gaits are violent: the trunk pitch standard deviation is
\SIrange{16}{57}{\degree}, and the servos deliver up to \SI{58}{W}. Patterns that are
symmetric about the sagittal plane (front--back, in-phase) have zero yaw by
construction; the best diagonal gaits drifted by \SIrange{2}{13}{\degree} in \SI{8}{s}.

\begin{figure}[t]
  \centering
  \includegraphics[width=\linewidth]{fig_gait_patterns}
  \caption{Toy robot, best gait per phase pattern after 150 rollouts in the 12-D space.
  Bars: mean of two seeds; dot and cross: seeds 1 and 2.}
  \label{fig:mj-gait-patterns}
\end{figure}

\paragraph{Duty cycle.}
Varying only the duty cycle of the best diagonal gait (optimum $d=0.54$) leaves the speed
between \SI{21.6}{cm/s} and \SI{22.4}{cm/s} over $d\in[0.35,0.65]$. It halves at the ends
of the range (\cref{fig:mj-duty}). The best front--back gait ($d=0.56$) is fragile: five of
the eleven duty values of the sweep lead to diverged rollouts. A robust optimum is preferable to a slightly
faster one on the edge of instability, and the diagonal pattern is the better candidate.

\begin{figure}[t]
  \centering
  \includegraphics[width=0.55\linewidth]{fig_duty}
  \caption{Speed versus power-stroke duty cycle for the best diagonal and front--back
  gaits, all other parameters fixed. Crosses: diverged rollouts. Dotted lines: optimised
  duty.}
  \label{fig:mj-duty}
\end{figure}

\subsection{The optimiser chooses lift-based thrust}
\label{sec:mj-liftdrag}
The flipper forces of the optimised gaits were decomposed along the relative flow. With
$\hat{\bm w}=-\bm v/|\bm v|$, the component $(\bm F\cdot\hat{\bm w})\hat{\bm w}$ is the
drag part and the remainder the lift part. Their cycle-averaged forward components are
shown in \cref{fig:mj-liftdrag}. In all six structured optima the forward thrust comes
entirely from the lift part (\SIrange{0.47}{1.77}{N} for four flippers), while the drag
part points backwards (\SIrange{-0.16}{-1.07}{N}). This is the signature of lift-based
propulsion in the sense of Fish \cite{Fish1996}. It is not how a drag-based paddler such
as the muskrat works, whose power stroke pushes water backwards \cite{Fish1984}. The
flippers sweep sideways and vertically as much as backwards (transverse path share
\numrange{0.38}{0.63}).

Two points qualify the result. First, the thrust-weighted incidence is
\SIrange{58}{71}{\degree}, far above stall. The ``lift'' here is the flow-normal
component of a flat-plate normal force, not attached-flow lift, which a Morison element
cannot represent. Second, the two unconverged full-space results do not show the pattern
(one is drag-dominated, the other produces almost no thrust). The result holds for the
good gaits, not for every gait.

The same behaviour had been observed when running the toy model: all four legs swam in
a lift-type manner. This was the reason for building BODY2's lift-type fluke leg, and
the reason for adding a foil element that represents attached-flow lift and stall
explicitly (\cref{sec:mj-foil}).

\begin{figure}[t]
  \centering
  \includegraphics[width=\linewidth]{fig_liftdrag}
  \caption{Left: cycle-mean forward force of the four flippers, split into the part
  normal to the relative flow (lift) and the part along it (drag), for the structured
  optima and the two full-space runs. Right: transverse share of the flipper path and
  thrust-weighted incidence (colours as in \cref{fig:mj-gait-patterns}; grey: full space).}
  \label{fig:mj-liftdrag}
\end{figure}

% ---------------------------------------------------------------------
\section{Case B: BODY2 with passive-pitch flukes}
\label{sec:mj-fluke}

\paragraph{Model.}
The BODY2 model (\cref{fig:mj-models}, right) consists of a $262\times116\times\SI{50}{mm}$
hull (\SI{0.45}{kg}) with a low ballast block (\SI{0.85}{kg}) that places the centre
of mass below the centre of buoyancy. Each leg has a fixed thigh at
$\gamma=\SI{-85}{\degree}$, an actuated shank of \SI{47}{mm} whose angle $q_r$ is the
firmware's rocker angle, and a fluke on a pin with a passive pitch joint
(\cref{tab:mj-params}). The four-bar that drives the shank is not simulated; its
inverse kinematics are used only for feasibility and for converting the result to servo
angles. The shank servos are modelled as position actuators with $k_p=\SI{3}{N.m/rad}$
and $|\tau|\le\SI{0.8}{N.m}$. The gait uses one harmonic per shank: the passive pitch
joint supplies the stroke asymmetry that the second harmonic had to supply on the toy
robot. As reference, the gait exported to the firmware from an earlier run is used:
$f=\SI{1.4}{Hz}$, $q_r=\SI{38.2}{\degree}\pm\SI{23.1}{\degree}$, diagonal.

\subsection{Phase pattern}
At the reference frequency and amplitude, the diagonal pattern is the fastest
(\SI{1.59}{cm/s}) and also by far the calmest: pitch standard deviation
\SI{0.12}{\degree} and heave \SI{0.3}{mm} (\cref{fig:mj-fluke-phase}). Front--back and
in-phase strokes make the hull nod (\SIrange{3.1}{3.3}{\degree}, 25 times more) and heave
(\SIrange{4.4}{5.4}{mm}). This is consistent with the nodding reported from videos of the
physical robot, whose hind legs currently stroke in phase. In the diagonal pattern the vertical fluke forces of the two
diagonals cancel in pitch. The wave pattern turns the robot (\SI{47}{\degree} in
\SI{8}{s}), and left--right swims backwards with \SI{17}{\degree} roll. All further BODY2
results use the diagonal pattern.

\begin{figure}[t]
  \centering
  \includegraphics[width=\linewidth]{fig_fluke_phase}
  \caption{BODY2 at the reference stroke ($f=\SI{1.4}{Hz}$, $q_r=\SI{38.2}{\degree}\pm\SI{23.1}{\degree}$):
  speed, pitch standard deviation and heave standard deviation for five phase patterns.}
  \label{fig:mj-fluke-phase}
\end{figure}

\subsection{Hinge stiffness}
\Cref{fig:mj-fluke-stiffness} varies the leaf-spring stiffness. With a rigid fluke the
robot does not move ($|\bar u|<\SI{0.05}{cm/s}$). The passive pitch is what produces
thrust. A very soft hinge lets the fluke swing onto its \SI{40}{\degree} stop, and
thrust falls again. The optimum lies at $k\approx\SI{2.5}{mN.m/rad}$ for
\SI{1.4}{Hz} (the current spring) and at $k\approx\SI{1}{mN.m/rad}$ for \SI{1.0}{Hz}. The
optimal stiffness rises with frequency. This is consistent with the hydrodynamic moment
on the fluke growing with the square of the stroke speed, so that a stiffer spring is
needed for the same deflection.
If the stroke frequency is changed, the spring should be changed with it.

\begin{figure}[t]
  \centering
  \includegraphics[width=\linewidth]{fig_fluke_stiffness}
  \caption{BODY2, diagonal reference stroke: speed (left) and maximum passive pitch
  (right) versus hinge stiffness at two stroke frequencies. Dotted lines: rigid fluke.}
  \label{fig:mj-fluke-stiffness}
\end{figure}

\subsection{Frequency and amplitude}
\Cref{fig:mj-fluke-map} maps speed over stroke frequency and shank amplitude. For
amplitudes above \SI{22.6}{\degree} the mean shank angle is raised so that the lower
extreme stays at the \SI{15}{\degree} clearance limit. Speed rises monotonically with
amplitude, so the best stroke lies on the clearance constraint. Over frequency it peaks
near \SI{1.8}{Hz}. The best grid point, $f=\SI{1.8}{Hz}$ and $q_r=\SI{43.5}{\degree}\pm\SI{28}{\degree}$,
scores \SI{2.85}{cm/s}, against \SI{1.59}{cm/s} for the reference stroke. Power grows
faster than speed: going from the reference stroke to the best grid point raises the
speed by a factor of 1.8 and the servo power by a factor of 2.3 (right panel).

\paragraph{Transient bias of the score.}
The score is the mean speed over the first \SI{8}{s} of stroking and includes the
acceleration from rest. The surge time constant of BODY2 is several seconds, so this
underestimates the cruising speed. Rerunning both strokes for \SI{30}{s} and fitting
the last \SI{10}{s} gives steady speeds of \SI{3.28}{cm/s} (reference, \SI{67}{mW}) and
\SI{4.99}{cm/s} (best grid point, \SI{168}{mW}, \SI{0.17}{BL/s}), i.e.\ 2.1 and 1.75 times
the scored values. The ranking is unchanged. Absolute speeds should be quoted from long
runs, and the scored interval should be lengthened or restricted to the last cycles in
future searches.

\begin{figure}[t]
  \centering
  \includegraphics[width=\linewidth]{fig_fluke_map}
  \caption{BODY2, diagonal pattern. Left: scored speed over frequency and shank amplitude;
  white contours: maximum passive pitch; star: reference firmware stroke; circle: best
  grid point. Above the dashed line the mean angle is raised to respect the clearance
  limit. Right: speed versus mean servo power for each amplitude.}
  \label{fig:mj-fluke-map}
\end{figure}

\subsection{What the fluke actually does}
\Cref{fig:mj-fluke-history} shows two cycles of the best grid stroke after \SI{6}{s}.
The fluke pitches passively by $\pm\SI{26}{\degree}$, lagging the shank. Its angle of
attack still stays near $\pm\SIrange{60}{80}{\degree}$ for most of the stroke. The
attached-flow branch of the foil model ($\sigma=1$) is active in only
\SI{1.1}{\percent} of the time, and the model is fully stalled in
\SI{98.2}{\percent}. The reason is the low advance ratio: the three-quarter-chord point
moves at up to \SI{0.36}{m/s} while the robot cruises at \SI{0.05}{m/s}, so the inflow
arrives at about \SI{82}{\degree} to the stroke path. A passive pitch of more than
\SI{60}{\degree} would be needed to bring $\alpha$ into the attached range, beyond the
\SI{40}{\degree} stop. Each fluke produces a mean thrust of \SI{3.5}{mN} but a mean upward
force of \SI{11}{mN}, three times larger. Flapping foils operate efficiently at moderate
angles of attack and higher advance ratios \cite{Anderson1998}; at BODY2's speeds the
fluke works as a feathering flat plate, and its thrust depends on the post-stall
normal-force coefficient $C_N$, the least certain parameter of the model.

\begin{figure}[t]
  \centering
  \includegraphics[width=\linewidth]{fig_fluke_history}
  \caption{BODY2, best grid stroke ($f=\SI{1.8}{Hz}$, $q_r=\SI{43.5}{\degree}\pm\SI{28}{\degree}$),
  front-left leg over two cycles: shank angle and passive fluke pitch, angle of attack at
  the three-quarter chord (grey band: attached flow), fluke force components, and trunk
  surge speed.}
  \label{fig:mj-fluke-history}
\end{figure}

\begin{figure}[t]
  \centering
  \shot{shot_fluke_filmstrip.png}{\linewidth}{Four frames over one stroke of the best grid gait\\ ($f=1.8$\,Hz, $q_r=43.5^\circ\pm28^\circ$, diagonal), side view at water level,\\ frames at $t/T=0,\,0.25,\,0.5,\,0.75$}
  \caption{BODY2 swimming with the best grid stroke. The diagonal leg pairs stroke in
  antiphase; the flukes feather passively on the pin.}
  \label{fig:mj-filmstrip}
\end{figure}

\subsection{CMA-ES on BODY2}
CMA-ES was run on the seven-dimensional BODY2 space (frequency, amplitude, offset, four
phases) with the objective of \cref{eq:mj-objective}, three seeds and 144 rollouts each
(\cref{tab:mj-fluke-cma}). With infeasible candidates scored $-1$, 25--37\,\% of all
candidates were infeasible, and the incumbents reached at most \SI{2.1}{cm/s}, below the
best grid point. A repair step that shifts the mean shank angle into the feasible
interval raised feasibility to 93--95\,\% and produced candidates up to \SI{3.0}{cm/s}.
The incumbents were still slow gaits, because at speeds of 1--3\,cm/s the yaw and heave
penalties are as large as the speed term: a gait with free phases that swims faster but
turns by a few degrees scores worse than one that barely moves. The step size grew
instead of shrinking in the repaired runs.

For BODY2 the useful sequence is therefore: fix the phases to the diagonal pattern,
which is justified independently by \cref{fig:mj-fluke-phase}; handle the clearance
constraint by repair; and rescale the penalties to the speed range of the robot. The
remaining space (frequency, amplitude, offset, and optionally hinge stiffness) is small
enough for the sweeps above to cover it directly.

\begin{table}[t]
\centering
\caption{CMA-ES on BODY2, 144 rollouts per run. ``Best'' is the incumbent under
\cref{eq:mj-objective}; ``max.'' is the fastest feasible candidate sampled.}
\label{tab:mj-fluke-cma}
\small
\begin{tabular}{@{}lcccc@{}}
\toprule
constraint handling & seed & feasible & best: speed [cm/s] & max.\ speed [cm/s] \\
\midrule
penalty ($-1$)      & 1 & 63\,\% & 2.11 & 2.14 \\
                    & 2 & 73\,\% & 0.02 & 1.27 \\
                    & 3 & 75\,\% & 0.17 & 1.33 \\
repair (shift mean) & 1 & 95\,\% & 1.30 & 2.28 \\
                    & 2 & 93\,\% & $-0.01$ & 3.03 \\
\midrule
grid, diagonal (\cref{fig:mj-fluke-map}) & -- & -- & 2.85 & -- \\
\bottomrule
\end{tabular}
\end{table}

% ---------------------------------------------------------------------
\section{Discussion and limitations}
\label{sec:mj-discussion}

\paragraph{What the model can and cannot say.}
The reduced-order model is useful for questions whose answer depends on symmetry and
kinematics: whether a parameterisation can produce thrust at all, which phase patterns
cancel pitch, whether a passive hinge is needed, and where constraints bind. These
conclusions hold for any reasonable coefficient values. Absolute speeds, the optimal
stiffness and the optimal frequency depend directly on uncalibrated coefficients and
should be treated as hypotheses for experiments.

\paragraph{Main modelling limitations.}
\begin{itemize}
  \item \emph{Quasi-steady forces.} Wake, vortex shedding and the delayed-stall lift of
    impulsively started plates are absent. The single-leg CFD in \chCFD{} found unsteady
    forces about three times larger than the quasi-steady estimate for the paddle leg.
    Since the fluke operates almost entirely post-stall, its force level hinges on $C_N$
    and must be calibrated against CFD, UVLM or a thrust measurement at matching
    geometry, frequency and stiffness.
  \item \emph{Isotropic added mass.} The mass trick adds the same inertia in all
    directions and independent of immersion. For thin paddles and flukes this
    overestimates the inertia along the plate and underestimates the reaction normal to
    it during stroke reversal.
  \item \emph{No free surface.} Waves, spray and the change of immersion within an
    element are not modelled beyond \cref{eq:mj-immersion}. Strokes that leave the
    water are penalised only through the loss of immersion.
  \item \emph{Unmodelled mechanism.} The BODY2 four-bar and servo dynamics are replaced by
    a position servo on the shank; backlash and velocity-dependent torque limits are
    missing.
  \item \emph{Transient score.} The \SI{8}{s} score underestimates cruising speed by a
    factor of about two for BODY2 (\cref{sec:mj-fluke}).
\end{itemize}

\paragraph{Calibration plan.}
In order of expected effect: (i) weigh the assembled robot with ballast and set the
still-water height of O2 to \SI{-11}{mm}; (ii) a coast-down of the \emph{complete}
robot with legs at rest posture, fitted as in \cref{fig:mj-checks}, which calibrates the
resistance including appendages; (iii) the mean thrust and passive pitch amplitude of a
single fluke at known frequency, amplitude and stiffness, used to fit $C_N$ and $\alpha_s$;
(iv) a free-swimming run of the reference stroke, as an independent validation (target:
speed within \SI{20}{\percent}).

% ---------------------------------------------------------------------
\section{Summary}
\label{sec:mj-summary}
A robot-agnostic reduced-order simulation of paddling was built in MuJoCo from Morison
elements for hull and links and a quasi-steady foil element for passively pitching
flukes. It evaluates a whole-robot gait in about one to two seconds. Verification found
and removed three protocol errors that had inflated earlier results. On the toy robot,
the analysis showed that a first-harmonic gait cannot produce thrust with rigid
flippers, that a structured 12-dimensional search outperforms the full 35-dimensional
one, and that the optimiser drives the flippers into lift-based rather than drag-based
propulsion. On BODY2, the diagonal pattern is the fastest and suppresses nodding by a
factor of 25, the passive hinge is necessary and its best stiffness depends on
frequency, and the best stroke lies on the fluke--hull clearance limit. At the robot's
low advance ratio the flukes work almost entirely post-stall. Calibrating the
post-stall normal force and the whole-robot resistance is therefore the next step before
any predicted speed can be compared with the physical robot.
